\chapter{初めての到着：チャンギ空港と第一印象}

\section{出発と到着}

今年の1月、シンガポールで開催されるAAA I国際会議に参加する機会を得た。これは数十年の歴史を持つ人工知能の国際学会だ。長い間、コンピュータ科学の分野では特に輝かしいプラットフォームとみなされていなかったが、近年の人工知能の急速な台頭とともに、この会議の重みも増してきた。

旅程は1月17日の夜に始まった。ニューヨークのジョン・F・ケネディ空港を出発し、アブダビで乗り継いでシンガポールへ向かった。十数時間の移動を経て、1月19日の朝についに目的地に到着した。

\section{チャンギ空港：効率の最初のレッスン}

飛行機はチャンギ空港に着陸した。最も印象に残ったのは、入国手続きのスムーズさだった。私はAPECレーンを利用した。アジア太平洋経済協力（APEC）は加盟経済体のビジネス渡航者向けに専用の優先通関カードを提供しており、保有者は通常の入国カウンターの列に並ばずに済む。パスポートをリーダーにかざすだけで、およそ10秒で通関が完了した。多くの国で経験する繁雑な入国手続きと比べると、この効率には驚かされた。飛行機は午前9時半に着陸し、10時前にはターミナルの外に出ていた。

これがシンガポールから受けた最初の授業だったかもしれない――効率とはここではスローガンではなく、日常の運営に組み込まれた制度的な習慣なのだ。

\section{市内へ：アンクルとこの都市の輪郭}

ターミナルを出た後、スマートフォンでGrabを呼んだ。これはシンガポールで最も普及しているライドシェアサービスで、すぐに運転手が応答した――華人の運転手だった。

シンガポールで初めて気づいた興味深い呼び方がある。運転手が華人の男性なら「Uncle（おじさん）」、女性なら「Auntie（おばさん）」と呼ぶのが一般的だ。この呼び方には家族的な親しみがあり、思わず微笑んでしまう。後で知ったのだが、この習慣は主に華人コミュニティの中で使われており、マレー系やインド系の運転手には通常使われないとのことだった。

その日は日曜日の朝で、道路は驚くほど静かだった。チャンギ空港はシンガポール島の東端に位置しており、車は沿岸高速道路を通ってほぼ渋滞なく市内中心部へと向かった。道中、アンクル運転手が窓の外の景色を熱心に説明してくれた。新興商業地区、天に伸びる金融センター、そして「モール」と呼ばれる巨大なショッピング複合施設の数々。

およそ25分後、シンガポールの沿岸核心市街地に入った。

\section{ランドマーク：都市の自己表現}

最初に目に飛び込んできたのは、巨大な円形の構造物――シンガポール・フライヤーだ。その輪郭は朝の光の中で際立っており、海辺に浮かぶ巨大な時計のようだった。

少し遠くには、三棟の高層ビルが並んで立ち、その頂上に船の船体を思わせるアーチ型の構造物が架かっていた。世界的に有名なマリーナ・ベイ・サンズだ。シンガポールに2か所しかない合法カジノのうちの一つがここにある。その存在自体が興味深い政策的選択を物語っている――厳格な法治で知られるこの国が、都市のスカイラインの中心にカジノを配置しているのだ。

さらに遠くには、海へと続く広大な沿岸公園が広がっていた。

\section{アンダーズホテル：都市の垂直的論理}

車は海辺を離れて市内中心部に入り、いくつかの通りを曲がった後、高層ビルの入口前に停まった。ここが今回の宿泊先、アンダーズホテルだ。

このホテルは独立した建物ではなく、大規模な複合開発の一部として組み込まれている。建物に入るとエレベーターが入口階から直接25階まで上がる――そこがホテルのロビーだ。つまり、25階より下は商業スペース、オフィス、その他の用途に充てられており、ホテルは25階から上にある。

このように空間を垂直に積み重ねる利用方法は、土地が極めて高価なシンガポールでは珍しくない。土地の制約に対するこの都市の一貫した対応論理を反映している――上へと成長するのだ。

\section{答えの出ない問い}

私たちはシンガポールに計10日間滞在した。ほとんどの時間は会議への参加と仕事に費やされたが、ホテルが市内中心部から近かったため、毎朝と毎夕、少しの時間を街歩きに充てることができた。

こうした一人の時間の中で、私は同じ問いを何度も心の中で繰り返していた。

\begin{center}
\textit{なぜ中国は改革開放の初期に、この広大な大国がシンガポールのような小さな国から学ぼうとしたのか。}
\end{center}

広く語り継がれる逸話を思い出す。李光耀が初めて北京を訪問した際、鄧小平は「中国は先進国になれるだろうか」と尋ねたという。

李光耀は後にこう回顧した。北京に残った人々は中国のエリートたちだったが、海を渡ってシンガポールにやってきた華人の多くは、社会の底辺にいた労働者だった。そのような貧しい出自の移民たちが、ほぼ天然資源のない小さな島の上に近代国家を築けたのならば、中国にも近代化を達成する理由が十分あるはずだ、と。

その意味で、シンガポールは単なる都市国家ではない――それは現在進行形の実験だ。大国が林立する世界秩序の中で、小国はいかにして自らの生き残りの道を見つけるのか。
